\documentclass[12pt,a4paper]{article}

\usepackage[utf8]{inputenc}
\usepackage[T1]{fontenc}
\usepackage{amsmath, amssymb}
\usepackage{geometry}
\usepackage{xcolor}
\usepackage{tcolorbox}
\usepackage{enumitem}
\usepackage{fancyhdr}
\usepackage{hyperref}
\usepackage{mathtools}

\geometry{margin=2.5cm}
\setlength{\headheight}{14pt}
\tcbuselibrary{most}

\definecolor{opcolor}{gray}{0.5}
% Annotation dans un alignat: encadre le commentaire en mode texte gris
\newcommand{\ann}[1]{\qquad\textcolor{opcolor}{\text{\small #1}}}

\newtcolorbox{noteimportante}{
  colback=yellow!10!white, colframe=orange!80!black,
  title=\textbf{Note importante}, fonttitle=\bfseries, boxrule=1pt, arc=4pt}
\newtcolorbox{methode}{
  colback=blue!5!white, colframe=blue!60!black,
  title=\textbf{Méthode}, fonttitle=\bfseries, boxrule=1pt, arc=4pt}
\newtcolorbox{rappel}{
  colback=green!5!white, colframe=green!50!black,
  title=\textbf{Rappel}, fonttitle=\bfseries, boxrule=1pt, arc=4pt}
\newtcolorbox{solbox}{
  colback=gray!8!white, colframe=black!60,
  title=\textbf{Solution}, fonttitle=\bfseries, boxrule=1pt, arc=4pt}

\pagestyle{fancy}
\fancyhf{}
\rhead{Systèmes Linéaires --- L1 Informatique}
\lhead{TD1 --- Correction détaillée}
\cfoot{\thepage}
\renewcommand{\headrulewidth}{0.4pt}
\hypersetup{hidelinks}

\begin{document}

\begin{center}
  {\LARGE \textbf{TD1 --- Systèmes d'équations linéaires}}\\[0.5em]
  {\large Algèbre et Programmation --- Correction détaillée}\\[0.3em]
  {\normalsize Licence 1 Informatique --- CERI --- Avignon Université --- 2026}
\end{center}
\vspace{0.5em}\hrule\vspace{1em}

%=====================================================================
\section*{Introduction et rappels}
%=====================================================================

\begin{rappel}
\textbf{Qu'est-ce qu'un système d'équations linéaires ?}

Un système linéaire est un ensemble d'équations où les inconnues ($x, y, z\ldots$) n'apparaissent qu'au premier degré. On cherche toutes les valeurs qui satisfont \textbf{simultanément} toutes les équations.

\medskip
\textbf{Les 3 opérations autorisées (qui ne changent pas les solutions) :}
\begin{itemize}[noitemsep]
  \item Échanger deux équations : $(E_i) \leftrightarrow (E_j)$
  \item Multiplier une équation par un réel $\lambda \neq 0$ : $(E_i) \leftarrow \lambda \cdot (E_i)$
  \item Ajouter un multiple d'une équation à une autre : $(E_j) \leftarrow (E_j) + \lambda \cdot (E_i)$
\end{itemize}

\medskip
\textbf{Les 3 cas possibles :}
\begin{itemize}[noitemsep]
  \item \textbf{Solution unique} : les droites/plans se coupent en un seul point
  \item \textbf{Aucune solution} : contradiction $0 = \alpha$ (avec $\alpha \neq 0$) après élimination
  \item \textbf{Infinité de solutions} : équation redondante $0 = 0$, il reste des variables libres
\end{itemize}
\end{rappel}

%=====================================================================
\section{Systèmes simples}
%=====================================================================

\textit{Résoudre les systèmes linéaires suivants.}

%---------------------------------------------------------------------
\subsection{Système (a)}
%---------------------------------------------------------------------

\[
\begin{cases}
  x + y + 2z = 3 & (1)\\
  x + 2y + z = 1 & (2)\\
  2x + y + z = 0 & (3)
\end{cases}
\]

\subsubsection*{$\triangleright$ Résolution ligne par ligne}

\textbf{Étape 1 :} Éliminer $x$ des équations (2) et (3).

\begin{alignat*}{4}
  x &{}+{}& y  &{}+{}& 2z &{}={}& 3   &\quad (1)\\
    &     & y  &{}-{}& z  &{}={}& -2  &\quad (2') \ann{$(2)-(1)$}\\
    &{}-{}& y  &{}-{}& 3z &{}={}& -6  &\quad (3') \ann{$(3) - 2\times(1)$}
\end{alignat*}

\textbf{Étape 2 :} Éliminer $y$ de (3').

\begin{alignat*}{4}
  x &{}+{}& y  &{}+{}& 2z &{}={}& 3   &\quad (1)\\
    &     & y  &{}-{}& z  &{}={}& -2  &\quad (2')\\
    &     &    &{}-{}& 4z &{}={}& -8  &\quad (3'') \ann{$(3') + (2')$}
\end{alignat*}

Le système est \textbf{triangulaire}. On résout de bas en haut.

\textbf{Étape 3 : substitution arrière.}

Depuis $(3'')$: $\quad -4z = -8 \implies \boxed{z=2}$

Depuis $(2')$: $\quad y - z = -2 \implies y = -2 + 2 = \boxed{0}$

Depuis $(1)$: $\quad x = 3 - y - 2z = 3 - 0 - 4 = \boxed{-1}$

\begin{solbox}
$\mathcal{S} = \{(-1,\ 0,\ 2)\}$ \hfill \textbf{Vérif. dans (3) :} $2(-1) + 0 + 2 = 0$ \checkmark
\end{solbox}

\subsubsection*{$\triangleright$ Méthode des matrices augmentées}

\begin{methode}
La matrice augmentée $[A \mid b]$ représente le système en n'écrivant que les coefficients sans les noms des variables. Chaque ligne = une équation. Les mêmes opérations s'appliquent sur les lignes.
\end{methode}

\[
\begin{pmatrix}
1 & 1 & 2 & \mid & 3 \\
1 & 2 & 1 & \mid & 1 \\
2 & 1 & 1 & \mid & 0
\end{pmatrix}
\xrightarrow{\substack{L_2 \leftarrow L_2 - L_1 \\ L_3 \leftarrow L_3 - 2L_1}}
\begin{pmatrix}
1 & 1 & 2 & \mid & 3 \\
0 & 1 & -1 & \mid & -2 \\
0 & -1 & -3 & \mid & -6
\end{pmatrix}
\xrightarrow{L_3 \leftarrow L_3 + L_2}
\begin{pmatrix}
1 & 1 & 2 & \mid & 3 \\
0 & 1 & -1 & \mid & -2 \\
0 & 0 & -4 & \mid & -8
\end{pmatrix}
\]

On retrouve $z=2$, $y=0$, $x=-1$.

%---------------------------------------------------------------------
\subsection{Système (b)}
%---------------------------------------------------------------------

\[
\begin{cases}
  x + 2z = 1 & (1)\\
  {-y} + z = 2 & (2)\\
  x - 2y = 1 & (3)
\end{cases}
\]

\begin{noteimportante}
Certains coefficients sont \textbf{absents} : cela signifie qu'ils valent zéro. Par exemple, $y$ n'apparaît pas dans (1), ce qui est équivalent à écrire $1\cdot x + \mathbf{0}\cdot y + 2\cdot z = 1$. Il faut en tenir compte lors des opérations sur les lignes.
\end{noteimportante}

\subsubsection*{$\triangleright$ Résolution ligne par ligne}

\textbf{Étape 1 :} Éliminer $x$ de (3).

\begin{alignat*}{4}
  x &      &    &{}+{}& 2z  &{}={}& 1   &\quad (1)\\
    &{}-{}& y  &{}+{}& z   &{}={}& 2   &\quad (2)\\
    &{}-{}& 2y &{}-{}& 2z  &{}={}& 0   &\quad (3') \ann{$(3)-(1)$}
\end{alignat*}

\textbf{Étape 2 :} Éliminer $y$ de (3').

\begin{alignat*}{4}
  x &      &    &{}+{}& 2z  &{}={}& 1   &\quad (1)\\
    &{}-{}& y  &{}+{}& z   &{}={}& 2   &\quad (2)\\
    &      &    &{}-{}& 4z  &{}={}& -4  &\quad (3'') \ann{$(3') - 2\times(2)$}
\end{alignat*}

\textbf{Étape 3 : substitution arrière.}

Depuis $(3'')$: $\quad -4z = -4 \implies \boxed{z=1}$

Depuis $(2)$: $\quad -y + z = 2 \implies -y = 2-1 = 1 \implies \boxed{y=-1}$

Depuis $(1)$: $\quad x + 2z = 1 \implies x = 1-2 = \boxed{-1}$

\begin{solbox}
$\mathcal{S} = \{(-1,\ -1,\ 1)\}$ \hfill \textbf{Vérif. dans (3) :} $(-1) - 2(-1) = 1$ \checkmark
\end{solbox}

\subsubsection*{$\triangleright$ Méthode des matrices augmentées}

\[
\begin{pmatrix}
1 & 0 & 2 & \mid & 1 \\
0 & -1 & 1 & \mid & 2 \\
1 & -2 & 0 & \mid & 1
\end{pmatrix}
\xrightarrow{L_3 \leftarrow L_3 - L_1}
\begin{pmatrix}
1 & 0 & 2 & \mid & 1 \\
0 & -1 & 1 & \mid & 2 \\
0 & -2 & -2 & \mid & 0
\end{pmatrix}
\xrightarrow{L_3 \leftarrow L_3 - 2L_2}
\begin{pmatrix}
1 & 0 & 2 & \mid & 1 \\
0 & -1 & 1 & \mid & 2 \\
0 & 0 & -4 & \mid & -4
\end{pmatrix}
\]

%=====================================================================
\section{Plus d'inconnues ou d'équations}
%=====================================================================

\textit{Résoudre les systèmes suivants.}

%---------------------------------------------------------------------
\subsection{Système (a) --- 2 équations, 4 inconnues}
%---------------------------------------------------------------------

\[
\begin{cases}
  x + y + z - 3t = 1 & (1)\\
  2x + y - z + t = 1 & (2)
\end{cases}
\]

\begin{methode}
Quand il y a \textbf{moins d'équations que d'inconnues}, il reste forcément des \textbf{variables libres} : des variables sans pivot qu'on peut choisir librement. On les nomme $\lambda, \mu, \ldots$ et on exprime les autres en fonction d'elles.
\end{methode}

\subsubsection*{$\triangleright$ Résolution ligne par ligne}

\textbf{Étape 1 :} Éliminer $x$.

\begin{alignat*}{5}
  x &{}+{}& y  &{}+{}& z   &{}-{}& 3t &{}={}& 1  &\quad (1)\\
    &{}-{}& y  &{}-{}& 3z  &{}+{}& 7t &{}={}& -1 &\quad (2') \ann{$(2)-2\times(1)$}
\end{alignat*}

Les variables $z$ et $t$ n'ont pas de pivot : elles sont \textbf{libres}. On pose $z = \lambda$ et $t = \mu$.

\textbf{Étape 2 :} Exprimer $x$ et $y$ en fonction de $\lambda$ et $\mu$.

Depuis $(2')$:
\[
-y = -1 + 3\lambda - 7\mu \implies y = 1 - 3\lambda + 7\mu
\]

Depuis $(1)$:
\[
x = 1 - y - z + 3t = 1 - (1-3\lambda+7\mu) - \lambda + 3\mu = 2\lambda - 4\mu
\]

\begin{solbox}
Infinité de solutions paramétrées par $\lambda, \mu \in \mathbb{R}$ :
\[
\mathcal{S} = \bigl\{(2\lambda-4\mu,\;\; 1-3\lambda+7\mu,\;\; \lambda,\;\; \mu) \mid \lambda, \mu \in \mathbb{R}\bigr\}
\]
Sous forme vectorielle :
\[
\begin{pmatrix} x \\ y \\ z \\ t \end{pmatrix}
= \begin{pmatrix} 0 \\ 1 \\ 0 \\ 0 \end{pmatrix}
+ \lambda \begin{pmatrix} 2 \\ -3 \\ 1 \\ 0 \end{pmatrix}
+ \mu \begin{pmatrix} -4 \\ 7 \\ 0 \\ 1 \end{pmatrix}
\]
\end{solbox}

\subsubsection*{$\triangleright$ Méthode des matrices augmentées}

\[
\begin{pmatrix}
1 & 1 & 1 & -3 & \mid & 1 \\
2 & 1 & -1 & 1 & \mid & 1
\end{pmatrix}
\xrightarrow{L_2 \leftarrow L_2 - 2L_1}
\begin{pmatrix}
1 & 1 & 1 & -3 & \mid & 1 \\
0 & -1 & -3 & 7 & \mid & -1
\end{pmatrix}
\]

Colonnes 3 et 4 sans pivot $\Rightarrow$ $z$ et $t$ libres.

%---------------------------------------------------------------------
\subsection{Système (b) --- 4 équations, 3 inconnues}
%---------------------------------------------------------------------

\[
\begin{cases}
  x + 2y - 3z = 4 & (1)\\
  x + 3y - z = 11 & (2)\\
  2x + 5y - 5z = 13 & (3)\\
  x + 4y + z = 18 & (4)
\end{cases}
\]

\begin{methode}
Quand il y a \textbf{plus d'équations que d'inconnues}, certaines peuvent être redondantes ($\to 0=0$, elles disparaissent) ou incompatibles ($\to 0=\alpha\neq 0$, aucune solution). L'élimination permet de distinguer les deux cas.
\end{methode}

\subsubsection*{$\triangleright$ Résolution ligne par ligne}

\textbf{Étape 1 :} Éliminer $x$ des équations (2), (3), (4).

\begin{alignat*}{4}
  x &{}+{}& 2y  &{}-{}& 3z &{}={}& 4   &\quad (1)\\
    &     & y   &{}+{}& 2z &{}={}& 7   &\quad (2') \ann{$(2)-(1)$}\\
    &     & y   &{}+{}& z  &{}={}& 5   &\quad (3') \ann{$(3)-2\times(1)$}\\
    &     & 2y  &{}+{}& 4z &{}={}& 14  &\quad (4') \ann{$(4)-(1)$}
\end{alignat*}

\textbf{Étape 2 :} Éliminer $y$ des équations (3') et (4').

\begin{alignat*}{4}
  x &{}+{}& 2y  &{}-{}& 3z &{}={}& 4   &\quad (1)\\
    &     & y   &{}+{}& 2z &{}={}& 7   &\quad (2')\\
    &     &     &{}-{}& z  &{}={}& -2  &\quad (3'') \ann{$(3')-(2')$}\\
    &     &     &     & 0  &{}={}& 0   &\quad (4'') \ann{$(4')-2\times(2')$}
\end{alignat*}

L'équation (4'') donne $0=0$ : elle est \textbf{redondante} et disparaît.

\textbf{Étape 3 : substitution arrière.}

Depuis $(3'')$: $\quad \boxed{z=2}$

Depuis $(2')$: $\quad y = 7 - 2\times 2 = \boxed{3}$

Depuis $(1)$: $\quad x = 4 - 2\times 3 + 3\times 2 = \boxed{4}$

\begin{solbox}
$\mathcal{S} = \{(4,\ 3,\ 2)\}$ \hfill (solution unique malgré 4 équations : la 4\textsuperscript{ème} était redondante)
\end{solbox}

\subsubsection*{$\triangleright$ Méthode des matrices augmentées}

\[
\begin{pmatrix}
1 & 2 & -3 & \mid & 4 \\
1 & 3 & -1 & \mid & 11 \\
2 & 5 & -5 & \mid & 13 \\
1 & 4 & 1 & \mid & 18
\end{pmatrix}
\xrightarrow{\substack{L_2-L_1 \\ L_3-2L_1 \\ L_4-L_1}}
\begin{pmatrix}
1 & 2 & -3 & \mid & 4 \\
0 & 1 & 2 & \mid & 7 \\
0 & 1 & 1 & \mid & 5 \\
0 & 2 & 4 & \mid & 14
\end{pmatrix}
\xrightarrow{\substack{L_3-L_2 \\ L_4-2L_2}}
\begin{pmatrix}
1 & 2 & -3 & \mid & 4 \\
0 & 1 & 2 & \mid & 7 \\
0 & 0 & -1 & \mid & -2 \\
0 & 0 & 0 & \mid & 0
\end{pmatrix}
\]

%=====================================================================
\section{Systèmes pièges --- deux systèmes proches}
%=====================================================================

\textit{Résoudre les deux systèmes suivants et commenter.}

\[
(S_1) \;
\begin{cases}
  x + 5y + 9z = 180 \\
  9x + 10y + 5z = 40 \\
  10x + 9y + z = -50
\end{cases}
\qquad
(S_2) \;
\begin{cases}
  x + 5y + 9z = 180 \\
  9x + 10y + 5z = 41 \\
  10x + 9y + z = -50
\end{cases}
\]

\subsection*{Résolution de $(S_1)$}

\subsubsection*{$\triangleright$ Résolution ligne par ligne}

\textbf{Étape 1 :} Éliminer $x$.

\begin{alignat*}{4}
  x &{}+{}&  5y  &{}+{}&  9z  &{}={}&  180   &\quad (1)\\
    &{}-{}& 35y  &{}-{}& 76z  &{}={}& -1580  &\quad (2') \ann{$(2)-9\times(1)$}\\
    &{}-{}& 41y  &{}-{}& 89z  &{}={}& -1850  &\quad (3') \ann{$(3)-10\times(1)$}
\end{alignat*}

\textbf{Étape 2 :} Éliminer $y$. Les coefficients de $y$ sont $-35$ et $-41$. On calcule $35\times(3') - 41\times(2')$ :

Coefficient de $z$ : $35\times(-89) - 41\times(-76) = -3115 + 3116 = 1$

Membre droit : $35\times(-1850) - 41\times(-1580) = -64750 + 64780 = 30$

\begin{alignat*}{4}
  x &{}+{}&  5y  &{}+{}& 9z &{}={}&  180   &\quad (1)\\
    &{}-{}& 35y  &{}-{}& 76z &{}={}& -1580  &\quad (2')\\
    &     &      &     &  z  &{}={}&  30    &\quad (3'') \ann{$35\times(3') - 41\times(2')$}
\end{alignat*}

\textbf{Étape 3 :}

$z = 30$

Depuis $(2')$: $-35y = -1580 + 76\times 30 = -1580 + 2280 = 700 \implies y = -20$

Depuis $(1)$: $x = 180 - 5\times(-20) - 9\times 30 = 180 + 100 - 270 = 10$

\begin{solbox}
$\mathcal{S}_1 = \{(10,\ -20,\ 30)\}$
\end{solbox}

\subsection*{Résolution de $(S_2)$ --- seul le membre droit de (2) change : $40 \to 41$}

\textbf{Étape 1 :} $(2') \leftarrow (2) - 9\times(1)$ donne maintenant $-35y - 76z = \mathbf{-1579}$

\textbf{Étape 2 :} $35\times(3') - 41\times(2')$ :

Membre droit : $35\times(-1850) - 41\times(-1579) = -64750 + 64739 = -11$

Donc $z = -11$.

\textbf{Étape 3 :}

$z = -11$

Depuis $(2')$: $-35y = -1579 + 76\times(-11) = -1579 - 836 = -2415 \implies y = 69$

Depuis $(1)$: $x = 180 - 5\times 69 - 9\times(-11) = 180 - 345 + 99 = -66$

\begin{solbox}
$\mathcal{S}_2 = \{(-66,\ 69,\ -11)\}$
\end{solbox}

\begin{noteimportante}
\textbf{Commentaire --- système mal conditionné :}

Un seul coefficient change de 40 à 41 ($< 1\%$ de variation), et pourtant la solution passe de $(10, -20, 30)$ à $(-66, 69, -11)$. C'est un changement énorme !

On dit que ce système est \textbf{mal conditionné} : une petite perturbation des données peut produire de très grandes variations dans la solution. C'est un problème crucial en informatique numérique, car les ordinateurs travaillent avec des nombres arrondis --- même de très petites erreurs d'arrondi peuvent totalement fausser le résultat.
\end{noteimportante}

%=====================================================================
\section{Étude paramétrique selon $m \in \mathbb{R}$}
%=====================================================================

\textit{Déterminer, selon la valeur de $m \in \mathbb{R}$, l'ensemble des solutions du système :}

\[
\begin{cases}
  x + y - z = 1 & (1)\\
  3x + y - z = 1 & (2)\\
  2x - 2y + 2z = m & (3)
\end{cases}
\]

\begin{methode}
Quand un paramètre apparaît, on fait l'élimination normalement. On aboutit à une ligne entièrement nulle à gauche : le membre droit est alors une expression en $m$. On discute :
\begin{itemize}[noitemsep]
  \item Membre droit $= 0$ : équation redondante $\Rightarrow$ variables libres, infinité de solutions
  \item Membre droit $\neq 0$ : contradiction $\Rightarrow$ aucune solution
\end{itemize}
\end{methode}

\subsubsection*{$\triangleright$ Résolution ligne par ligne}

\textbf{Étape 1 :} Éliminer $x$.

\begin{alignat*}{4}
  x &{}+{}&  y  &{}-{}& z  &{}={}&  1    &\quad (1)\\
    &{}-{}& 2y  &{}+{}& 2z &{}={}& -2    &\quad (2') \ann{$(2)-3\times(1)$}\\
    &{}-{}& 4y  &{}+{}& 4z &{}={}& m-2   &\quad (3') \ann{$(3)-2\times(1)$}
\end{alignat*}

\textbf{Étape 2 :} Éliminer $y$ : $(3') - 2\times(2')$.

Membre gauche : $-4y+4z - 2(-2y+2z) = 0$

Membre droit : $(m-2) - 2\times(-2) = m + 2$

\begin{alignat*}{4}
  x &{}+{}&  y  &{}-{}& z  &{}={}& 1    &\quad (1)\\
    &{}-{}& 2y  &{}+{}& 2z &{}={}& -2   &\quad (2')\\
    &     &  0  &     &    &{}={}& m+2  &\quad (3'') \ann{$(3')-2\times(2')$}
\end{alignat*}

\paragraph{Cas $m \neq -2$ :} $0 = m+2 \neq 0$ : contradiction. $\mathcal{S} = \emptyset$.

\paragraph{Cas $m = -2$ :} $0 = 0$ : équation redondante. $z$ est libre ($z = \lambda$).

Depuis $(2')$: $-2y + 2\lambda = -2 \implies y = \lambda + 1$

Depuis $(1)$: $x = 1 - y + z = 1 - (\lambda+1) + \lambda = 0$

\begin{solbox}
\begin{itemize}[noitemsep]
  \item Si $m \neq -2$ : $\mathcal{S} = \emptyset$ (aucune solution)
  \item Si $m = -2$ : infinité de solutions
  \[
    \mathcal{S} = \left\{\begin{pmatrix} 0 \\ 1 \\ 0 \end{pmatrix}
    + \lambda \begin{pmatrix} 0 \\ 1 \\ 1 \end{pmatrix}
    \;\middle|\; \lambda \in \mathbb{R}\right\}
  \]
\end{itemize}
\end{solbox}

\subsubsection*{$\triangleright$ Méthode des matrices augmentées}

\[
\begin{pmatrix}
1 & 1 & -1 & \mid & 1 \\
3 & 1 & -1 & \mid & 1 \\
2 & -2 & 2 & \mid & m
\end{pmatrix}
\xrightarrow{\substack{L_2-3L_1 \\ L_3-2L_1}}
\begin{pmatrix}
1 & 1 & -1 & \mid & 1 \\
0 & -2 & 2 & \mid & -2 \\
0 & -4 & 4 & \mid & m-2
\end{pmatrix}
\xrightarrow{L_3-2L_2}
\begin{pmatrix}
1 & 1 & -1 & \mid & 1 \\
0 & -2 & 2 & \mid & -2 \\
0 & 0 & 0 & \mid & m+2
\end{pmatrix}
\]

%=====================================================================
\section{Résolution avec paramètre $m$ --- système $2\times2$}
%=====================================================================

\textit{Soit $m$ un réel. Résoudre le système et donner une interprétation géométrique.}

\[
\begin{cases}
  x + my = -3 & (1)\\
  mx + 4y = 6 & (2)
\end{cases}
\]

\subsubsection*{$\triangleright$ Résolution ligne par ligne}

\textbf{Étape 1 :} Éliminer $x$ avec $(2) - m\times(1)$ :

Membre gauche : $mx - mx + 4y - m^2 y = (4-m^2)y$

Membre droit : $6 - m\times(-3) = 6+3m$

\begin{alignat*}{3}
  x &{}+{}&    my       &{}={}& -3     &\quad (1)\\
    &     & (4-m^2)\,y  &{}={}& 6+3m   &\quad (2') \ann{$(2) - m\times(1)$}
\end{alignat*}

Le coefficient de $y$ vaut $4-m^2 = (2-m)(2+m)$.

\paragraph{Cas $m \neq 2$ et $m \neq -2$ :}
\[
y = \frac{6+3m}{4-m^2} = \frac{3(2+m)}{(2-m)(2+m)} = \frac{3}{2-m}
\qquad
x = -3 - my = -3 - \frac{3m}{2-m} = \frac{-6}{2-m}
\]

\paragraph{Cas $m = 2$ :} $(2')$ donne $0\cdot y = 12$, soit $0=12$ : \textbf{contradiction, aucune solution}.

\paragraph{Cas $m = -2$ :} $(2')$ donne $0\cdot y = 0$, soit $0=0$ : équation redondante, $y = \lambda$ libre.

Depuis $(1)$: $x = -3 - (-2)\lambda = -3 + 2\lambda$

\begin{solbox}
\begin{itemize}[noitemsep]
  \item $m \neq \pm 2$ : \textbf{solution unique} $\displaystyle\quad x = \frac{-6}{2-m},\quad y = \frac{3}{2-m}$
  \item $m = 2$ : \textbf{aucune solution} (droites parallèles)
  \item $m = -2$ : \textbf{infinité de solutions} $\quad \mathcal{S} = \{(-3+2\lambda,\ \lambda) \mid \lambda \in \mathbb{R}\}$ (droites confondues)
\end{itemize}

\smallskip
\textbf{Interprétation :} chaque équation représente une droite dans le plan $(x, y)$. Les 3 cas correspondent aux 3 configurations possibles de deux droites : sécantes, parallèles, ou confondues.
\end{solbox}

\subsubsection*{$\triangleright$ Méthode des matrices augmentées}

\[
\begin{pmatrix}
1 & m & \mid & -3 \\
m & 4 & \mid & 6
\end{pmatrix}
\xrightarrow{L_2 - m \cdot L_1}
\begin{pmatrix}
1 & m & \mid & -3 \\
0 & 4-m^2 & \mid & 6+3m
\end{pmatrix}
\]

Même discussion selon $4 - m^2$.

%=====================================================================
\section{Paramètres $a, b \in \mathbb{R}$}
%=====================================================================

\textit{Déterminer suivant les valeurs de $a, b \in \mathbb{R}$ l'ensemble des solutions du système :}

\[
\begin{cases}
  ax + y = b & (1)\\
  x + ay = b & (2)
\end{cases}
\]

\subsubsection*{$\triangleright$ Résolution ligne par ligne}

On échange (1) et (2) pour avoir un coefficient $1$ devant $x$ en tête :

\begin{alignat*}{3}
  x &{}+{}& ay         &{}={}& b       &\quad (1) \;\,\ann{ancienne (2)}\\
  ax &{}+{}& y          &{}={}& b       &\quad (2) \;\,\ann{ancienne (1)}
\end{alignat*}

\textbf{Étape 1 :} Éliminer $x$ avec $(2) - a\times(1)$ :

\begin{alignat*}{3}
  x &{}+{}& ay          &{}={}& b       &\quad (1)\\
    &     & (1-a^2)\,y  &{}={}& b(1-a)  &\quad (2') \ann{$(2)-a\times(1)$}
\end{alignat*}

On factorise : $1-a^2 = (1-a)(1+a)$.

\paragraph{Cas $a \neq 1$ et $a \neq -1$ :}
\[
y = \frac{b(1-a)}{(1-a)(1+a)} = \frac{b}{1+a}
\qquad
x = b - ay = b - \frac{ab}{1+a} = \frac{b}{1+a}
\]

\paragraph{Cas $a = 1$ :} $(2')$ donne $0\cdot y = 0$ : $0=0$ pour tout $b$. Équation redondante.

$y=\lambda$ libre. Depuis $(1)$: $x = b - \lambda$. \textbf{Infinité de solutions pour tout $b$.}

\paragraph{Cas $a = -1$ :} $(2')$ donne $0\cdot y = 2b$, soit $0 = 2b$.
\begin{itemize}[noitemsep]
  \item $b \neq 0$ : $0 = 2b \neq 0$ : contradiction, \textbf{aucune solution}.
  \item $b = 0$ : $0 = 0$, équation redondante. $y = \lambda$ libre. Depuis $(1)$: $x = -\lambda$.
\end{itemize}

\begin{solbox}
\begin{itemize}
  \item $a \neq \pm 1$, $b$ quelconque : \textbf{solution unique} $\displaystyle\quad x = y = \frac{b}{1+a}$
  \item $a = 1$, $b$ quelconque : \textbf{infinité de solutions} $\quad \mathcal{S} = \{(b-\lambda,\ \lambda) \mid \lambda \in \mathbb{R}\}$
  \item $a = -1$, $b \neq 0$ : \textbf{aucune solution}
  \item $a = -1$, $b = 0$ : \textbf{infinité de solutions} $\quad \mathcal{S} = \{(-\lambda,\ \lambda) \mid \lambda \in \mathbb{R}\}$
\end{itemize}
\end{solbox}

\subsubsection*{$\triangleright$ Méthode des matrices augmentées}

\[
\begin{pmatrix}
a & 1 & \mid & b \\
1 & a & \mid & b
\end{pmatrix}
\xrightarrow{L_1 \leftrightarrow L_2}
\begin{pmatrix}
1 & a & \mid & b \\
a & 1 & \mid & b
\end{pmatrix}
\xrightarrow{L_2 - a \cdot L_1}
\begin{pmatrix}
1 & a & \mid & b \\
0 & 1-a^2 & \mid & b(1-a)
\end{pmatrix}
\]

Même discussion selon $1-a^2$.

%=====================================================================
\section*{Conseils pour les examens}
%=====================================================================

\begin{noteimportante}
\textbf{Méthodologie à suivre à chaque exercice :}
\begin{enumerate}
  \item Numéroter chaque équation $(1), (2), (3)\ldots$
  \item Choisir le premier coefficient non nul d'une colonne (le \textbf{pivot}), et éliminer cette variable dans toutes les autres équations en annotant chaque opération : ex. \textcolor{opcolor}{$(2) - 3\times(1)$}
  \item Répéter colonne par colonne jusqu'au système triangulaire.
  \item Si un paramètre est présent : \textbf{discuter} dès qu'un pivot peut s'annuler.
  \item Résoudre de \textbf{bas en haut} par substitution.
  \item \textbf{Vérifier} dans au moins une équation originale.
\end{enumerate}

\medskip
\textbf{Signaux à reconnaître pendant l'élimination :}
\begin{itemize}[noitemsep]
  \item Ligne $[\;0\;\cdots\;0\;\mid\;0\;]$ $\Rightarrow$ équation redondante --- variable libre possible
  \item Ligne $[\;0\;\cdots\;0\;\mid\;\alpha\;]$ avec $\alpha \neq 0$ $\Rightarrow$ contradiction --- aucune solution
  \item Colonne sans pivot $\Rightarrow$ variable libre --- poser $= \lambda$ (ou $\mu$, etc.)
\end{itemize}
\end{noteimportante}


\begin{document}

\begin{center}
  {\LARGE \textbf{TD1 --- Systèmes d'équations linéaires}}\\[0.5em]
  {\large Algèbre et Programmation}\\[0.3em]
  {\normalsize Licence 1 Informatique --- CERI --- Avignon Université --- 2026}
\end{center}
\vspace{0.5em}\hrule\vspace{1em}

\begin{rappel}
\textbf{Opérations autorisées (ne changent pas les solutions) :}
\begin{itemize}[noitemsep]
  \item Échanger deux équations : $(E_i) \leftrightarrow (E_j)$
  \item Multiplier une équation par un réel $\lambda \neq 0$ : $(E_i) \leftarrow \lambda \cdot (E_i)$
  \item Ajouter un multiple d'une équation à une autre : $(E_j) \leftarrow (E_j) + \lambda \cdot (E_i)$
\end{itemize}
\end{rappel}

%=====================================================================
\section{Systèmes simples}
%=====================================================================

Résoudre les systèmes linéaires suivants :

\[
(a)\quad
\begin{cases}
  x + y + 2z = 3 \\
  x + 2y + z = 1 \\
  2x + y + z = 0
\end{cases}
\qquad\qquad
(b)\quad
\begin{cases}
  x + 2z = 1 \\
  {-y} + z = 2 \\
  x - 2y = 1
\end{cases}
\]

\vspace{6cm}

%=====================================================================
\section{Plus d'inconnues ou d'équations}
%=====================================================================

Résoudre les systèmes suivants :

\[
(a)\quad
\begin{cases}
  x + y + z - 3t = 1 \\
  2x + y - z + t = 1
\end{cases}
\qquad\qquad
(b)\quad
\begin{cases}
  x + 2y - 3z = 4 \\
  x + 3y - z = 11 \\
  2x + 5y - 5z = 13 \\
  x + 4y + z = 18
\end{cases}
\]

\vspace{8cm}

%=====================================================================
\section{Systèmes pièges --- deux systèmes proches}
%=====================================================================

Résoudre les deux systèmes suivants. Qu'en pensez-vous ?

\[
(S_1)\quad
\begin{cases}
  x + 5y + 9z = 180 \\
  9x + 10y + 5z = 40 \\
  10x + 9y + z = -50
\end{cases}
\qquad\qquad
(S_2)\quad
\begin{cases}
  x + 5y + 9z = 180 \\
  9x + 10y + 5z = 41 \\
  10x + 9y + z = -50
\end{cases}
\]

\newpage

%=====================================================================
\section{Étude paramétrique selon $m \in \mathbb{R}$}
%=====================================================================

Déterminer, selon la valeur de $m \in \mathbb{R}$ en utilisant l'algorithme de Gauss, l'ensemble des solutions du système :

\[
\begin{cases}
  x + y - z = 1 \\
  3x + y - z = 1 \\
  2x - 2y + 2z = m
\end{cases}
\]

\vspace{8cm}

%=====================================================================
\section{Résolution avec paramètre $m$ --- système $2\times2$}
%=====================================================================

Soit $m$ un réel. Résoudre le système suivant :

\[
\begin{cases}
  x + my = -3 \\
  mx + 4y = 6
\end{cases}
\]

Quelle interprétation géométrique du résultat faites-vous ?

\vspace{8cm}

%=====================================================================
\section{Paramètres $a, b \in \mathbb{R}$}
%=====================================================================

Déterminer suivant les valeurs des paramètres $a, b \in \mathbb{R}$ l'ensemble des solutions du système :

\[
\begin{cases}
  ax + y = b \\
  x + ay = b
\end{cases}
\]

\end{document}

